\uuid{pK4m}
\titre{Propagation forward dans un RNN simple}
\niveau{L3}
\module{Analyse de données}
\chapitre{Réseaux de neurones}
\sousChapitre{Réseaux récurrents, LSTM et GRU}
\theme{État caché, activation tanh, séquence temporelle}
\auteur{Maxime Nguyen}
\datecreate{2026-05-13}
\organisation{AMSCC}
\difficulte{2}

\contenu{
	\texte{
		On considère un RNN simple à un seul état caché, défini par les équations :
		\[
		h_t = \tanh(W_h h_{t-1} + W_x x_t + b_h), \qquad y_t = W_y h_t + b_y
		\]
		avec les dimensions suivantes : entrée $x_t \in \mathbb{R}$, état caché $h_t \in \mathbb{R}^2$, sortie $y_t \in \mathbb{R}$.
		
		Les paramètres du réseau sont :
		\[
		W_x = \begin{pmatrix} 1 \\ -1 \end{pmatrix}, \quad
		W_h = \begin{pmatrix} 0.5 & 0 \\ 0 & 0.5 \end{pmatrix}, \quad
		b_h = \begin{pmatrix} 0 \\ 0 \end{pmatrix}, \quad
		W_y = \begin{pmatrix} 1 & 1 \end{pmatrix}, \quad
		b_y = 0
		\]
		L'état caché initial est $h_0 = \begin{pmatrix} 0 \\ 0 \end{pmatrix}$.
		On traite la séquence d'entrée $x_1 = 1$, $x_2 = 0$, $x_3 = -1$.
	}
	\begin{enumerate}
		\item
		\question{Calculer $h_1$, $h_2$, $h_3$ en appliquant la formule de récurrence pas à pas. On rappelle que $\tanh$ s'applique composante par composante.}
		\reponse{
			\textbf{Pas $t=1$.}
			\[
			h_1 = \tanh\!\left(W_h h_0 + W_x x_1\right)
			= \tanh\!\left(\begin{pmatrix}0\\0\end{pmatrix} + \begin{pmatrix}1\\-1\end{pmatrix}\right)
			= \begin{pmatrix}\tanh(1)\\\tanh(-1)\end{pmatrix}
			\approx \begin{pmatrix}0.762\\-0.762\end{pmatrix}
			\]
			\textbf{Pas $t=2$.}
			\[
			W_h h_1 = \begin{pmatrix}0.5 & 0\\0 & 0.5\end{pmatrix}\begin{pmatrix}\tanh(1)\\-\tanh(1)\end{pmatrix}
			= \begin{pmatrix}0.5\tanh(1)\\-0.5\tanh(1)\end{pmatrix}
			\approx \begin{pmatrix}0.381\\-0.381\end{pmatrix}, \qquad
			W_x x_2 = \begin{pmatrix}1\\-1\end{pmatrix}(0) = \begin{pmatrix}0\\0\end{pmatrix}
			\]
			\[
			h_2 = \tanh\!\left(\begin{pmatrix}0.381\\-0.381\end{pmatrix}\right)
			= \begin{pmatrix}\tanh(0.381)\\\tanh(-0.381)\end{pmatrix}
			\approx \begin{pmatrix}0.365\\-0.365\end{pmatrix}
			\]
			\textbf{Pas $t=3$.}
			\[
			W_h h_2 = \begin{pmatrix}0.5\tanh(0.381)\\-0.5\tanh(0.381)\end{pmatrix}
			\approx \begin{pmatrix}0.183\\-0.183\end{pmatrix}, \qquad
			W_x x_3 = \begin{pmatrix}1\\-1\end{pmatrix}(-1) = \begin{pmatrix}-1\\1\end{pmatrix}
			\]
			\[
			h_3 = \tanh\!\left(\begin{pmatrix}0.183\\-0.183\end{pmatrix} + \begin{pmatrix}-1\\1\end{pmatrix}\right)
			= \begin{pmatrix}\tanh(-0.817)\\\tanh(0.817)\end{pmatrix}
			\approx \begin{pmatrix}-0.674\\0.674\end{pmatrix}
			\]
		}
		
		\item
		\question{En déduire les sorties $y_1$, $y_2$, $y_3$.}
		\reponse{
			On applique $y_t = W_y h_t + b_y$. Puisque $W_y = \begin{pmatrix}1 & 1\end{pmatrix}$ et $b_y = 0$ :
			\[
			y_t = \begin{pmatrix}1 & 1\end{pmatrix}\begin{pmatrix}h_t^{(1)}\\h_t^{(2)}\end{pmatrix} = h_t^{(1)} + h_t^{(2)}
			\]
			\[
			y_1 = \tanh(1) + \tanh(-1) \approx 0.762 + (-0.762) = 0, \quad
			y_2 = \tanh(0.381) + \tanh(-0.381) \approx 0, \quad
			y_3 = \tanh(-0.817) + \tanh(0.817) \approx 0.
			\]
			La sortie est identiquement nulle car les deux composantes de $h_t$ sont systématiquement opposées, propriété induite par la structure antisymétrique de $W_x$ et la parité impaire de $\tanh$.
		}
		
		\item
		\question{Décrire qualitativement ce que fait ce réseau : comment l'état caché $h_t$ évolue-t-il lorsque l'entrée change de signe ?}
		\reponse{
			Les deux composantes de $h_t$ sont en permanence opposées ($h_t^{(2)} = -h_t^{(1)}$), pour trois raisons conjointes :
			\begin{itemize}
				\item $W_x$ injecte $x_t$ dans la première composante et $-x_t$ dans la seconde ;
				\item $W_h = 0.5\,I_2$ atténue l'état précédent de façon identique sur les deux composantes, préservant leur relation d'opposition ;
				\item $\tanh$ est une fonction impaire, donc $\tanh(-u) = -\tanh(u)$ conserve l'opposition.
			\end{itemize}
			Lorsque $x_t$ change de signe, les composantes s'inversent progressivement. Le réseau encode donc la polarité courante de l'entrée, atténuée par l'historique. La sortie $y_t$ étant toujours nulle, l'information pertinente réside uniquement dans l'état caché $h_t$.
		}
		
		\item
		\question{(\textit{Bonus}) Que se passerait-il si on remplaçait $\tanh$ par la fonction identité ? Réécrire alors $h_t$ comme une combinaison linéaire des entrées passées $x_1, \ldots, x_t$.}
		\indication{Développer la récurrence $h_t = W_h h_{t-1} + W_x x_t$ et itérer.}
		\reponse{
			Avec l'activation identité, la récurrence devient $h_t = W_h h_{t-1} + W_x x_t$. Par développement itératif (avec $h_0 = 0$) :
			\[
			h_t = W_h^{t-1}W_x x_1 + W_h^{t-2}W_x x_2 + \cdots + W_h^0 W_x x_t = \sum_{k=1}^{t} W_h^{t-k} W_x x_k.
			\]
			Puisque $W_h = 0.5\,I_2$, on a $W_h^{t-k} = (0.5)^{t-k} I_2$. En substituant :
			\[
			h_t = \sum_{k=1}^{t} (0.5)^{t-k} I_2\, W_x x_k = W_x \left(\sum_{k=1}^{t} (0.5)^{t-k} x_k\right).
			\]
			Le vecteur $h_t$ est proportionnel à $W_x = \begin{pmatrix}1\\-1\end{pmatrix}$, le facteur de proportionnalité étant une moyenne exponentielle pondérée des entrées passées, avec des poids $(0.5)^{t-k}$ décroissant géométriquement vers le passé.
		}
	\end{enumerate}
}